\chapter{Teorema Principal}

\section{Formulación}

\begin{theorem}[Equivalencia QCAL-Riemann]
La Hipótesis de Riemann es lógicamente equivalente a la QCAL-Hipótesis:
\[
\text{RH} \;\longleftrightarrow\; \forall n \in \mathbb{Z},\; \operatorname{Re}(\lambda_n) = \frac12
\]
donde $\lambda_n$ son los autovalores del operador de Dirac adélico
$\overline{\mathfrak{D}}$.
\end{theorem}

\begin{proof}[Idea de la demostración]
El isomorfismo espectral entre:
\begin{itemize}
  \item El espectro de $\overline{\mathfrak{D}}$ actuando sobre $L^2(\mathbb{A})$
  \item Los ceros no triviales de $\zeta(s)$
\end{itemize}
se construye mediante el operador de Montgomery-Dyson. La correlación de pares
sigue la estadística GUE (Gaussian Unitary Ensemble), y el espaciamiento medio
entre niveles es unitario.
\end{proof}

\section{Certificación del Sistema}

El sistema QCAL está certificado como completo y operativo:

\begin{center}
\begin{tabular}{ll}
\toprule
\textbf{Componente} & \textbf{Estado} \\
\midrule
Lean 4 Cuántico - Formalizado & $\checkmark$ \\
Qubits Tipológicos QCAL - Definidos & $\checkmark$ \\
Operadores QCAL - Implementados & $\checkmark$ \\
Ecuaciones Operativas - Establecidas & $\checkmark$ \\
Flujo de Renormalización - Verificado & $\checkmark$ \\
Sistema Completo - Validado & $\checkmark$ \\
\bottomrule
\end{tabular}
\end{center}
\end{document}
